\documentclass[a4paper,12pt]{article}
\usepackage[utf8]{inputenc}
\usepackage[T1]{fontenc}
\usepackage[italian]{babel}
\usepackage{geometry}
\usepackage{lmodern}
\usepackage{amsmath,amssymb}
\usepackage{enumitem}
\usepackage{hyperref}
\usepackage{listings}
\usepackage{xcolor}
\usepackage{graphicx}
\usepackage{pgfplots}
\usepackage{fancyhdr}
\usepackage{titlesec}
\usepackage{graphicx}
\titleformat{\section}
  {\normalfont\Large\bfseries\centering}  % Format: font size, bold, centered
  {\thesection}{1em}{} 
\pagestyle{fancy}
\pgfplotsset{compat=1.18}
\fancyhead{}
\fancyfoot{}
\fancyhead[L]{Gabriele Perna}          % Left side
\fancyhead[C]{Programmazione e algoritmica}     % Center
\fancyhead[R]{2025/2026}             % Right side
\fancyfoot[C]{\thepage}           % Page number in the center of footer
\let\oldsection\section
\renewcommand{\section}{\clearpage\oldsection}
\geometry{margin=2.5cm}

\title{Esercitazione 02-11-2025}
\author{}
\date{}

\begin{document}

\maketitle
\tableofcontents

\section{Esercizio 1}
\textbf{Trova un minimo o massimo in un array.}\\
\begin{itemize}
    \item Array ordinato = $\theta(1)$
    \item Array non ordinato = $\theta(n)$ (caso pessimo)
\end{itemize}
\textbf{Trova un elemento in un array.}
\begin{itemize}
    \item Array ordinato = $\theta(\log_2n)$
    \begin{enumerate}
        \item Si prende un elemento mediano e si analizza.
        \item Se è uguale si è trovato l'elemento (caso ottimo), altrimenti si prosegue
        \item Se l'elemento è maggiore di quello da trovare allora si analizzano gli elementi minori e viceversa.
        \item Ripetere, prendendo un elemento mediano al nuovo subarray.
    \end{enumerate}
    \item Array non ordinato = $\theta(n)$ (caso pessimo)
\end{itemize}


\section{Esercizio 2}

\section{Esercizio 3}

\section{Esercizio 4}

\section{Esercizio 5}

\section{Esercizio 6}

\section{Esercizio 7}

\section{Esercizio 8}

\end{document}